% !TeX encoding = UTF-8
\documentclass[variant=guia,cover=institutional,mode=screen]{ua-engineering-v6-article}
% !TeX encoding = UTF-8
\UASetUniversity{Universidad Austral}
\UASetFaculty{Facultad de Ingeniería}
\UASetDepartment{Departamento de Computación}
\UASetCampus{Pilar}
\UASetCareer{Ingeniería Informática}
\UASetCommission{Comisión A}
\UASetProfessor{Equipo docente}
\UASetSemester{Segundo cuatrimestre 2026}
\UASetCourse{Sistemas Distribuidos}
\UASetDate{2026}


\UASetClassNumber{11}
\UASetClassTitle{Resiliencia distribuida}
\UASetDocKind{Guía resuelta}

\title{Clase 11 --- Guía resuelta}
\author{\UAProfessor}
\date{\UADate}

\begin{document}

\section{Criterio de lectura}
Las resoluciones son modelos de razonamiento. Deben verificarse contra los supuestos del caso y no memorizarse como recetas independientes del dominio.

\section{Ejercicios y resoluciones completas}
\begin{uaexercise}
Explique por qué una dependencia lenta puede ser más peligrosa que una dependencia completamente caída.
\end{uaexercise}
\begin{uaanswer}
La idempotencia separa intentos físicos de una operación lógica. Se usa una clave estable, se persiste el resultado junto con el efecto y los reintentos devuelven ese mismo resultado. El costo es estado adicional, política de expiración y coordinación con el almacenamiento.

El argumento debe identificar la hipótesis, construir la cadena lógica o el contraejemplo y concluir exactamente qué propiedad queda demostrada. Una intuición sin secuencia verificable no es suficiente.

La evidencia apropiada sería curvas de carga, colas, breaker state y degradación del camino crítico. También debe indicarse una condición que refute la elección o haga preferible una solución más simple.
\end{uaanswer}

\begin{uaexercise}
Construya una cadena causal donde una degradación local termine generando una caída visible en el borde del sistema.
\end{uaexercise}
\begin{uaanswer}
La respuesta debe situarse en retry budgets, backoff, jitter, circuit breakers, backpressure, bulkheads y graceful degradation. Se comienza declarando el supuesto decisivo y el comportamiento observable; luego se recorre una ejecución nominal y otra bajo dependencia lenta, retry storm, cola sin límite, pool agotado y cascada. El mecanismo elegido debe vincularse con una garantía concreta, evitando afirmar propiedades fuera de su alcance.

Una solución completa organiza la decisión con P-F-G-S-T: define el problema, enumera fallas, fija la garantía, presenta el protocolo y reconoce el trade-off. Debe incluir estados intermedios y qué ve el usuario si la operación no termina.

La evidencia apropiada sería curvas de carga, colas, breaker state y degradación del camino crítico. También debe indicarse una condición que refute la elección o haga preferible una solución más simple.
\end{uaanswer}

\begin{uaexercise}
Explique por qué la resiliencia distribuida no consiste en “hacer más retries”.
\end{uaexercise}
\begin{uaanswer}
Porque el retry solo ayuda si la causa es transitoria. Si la causa es saturación o degradación persistente, reintentar agrega más trabajo sobre el componente degradado y empeora la situación. En esos casos, el retry amplifica la falla en vez de tolerarla.
\end{uaanswer}

\begin{uaexercise}
Explique la diferencia entre tolerar una falla y amplificarla.
\end{uaexercise}
\begin{uaanswer}
La idempotencia separa intentos físicos de una operación lógica. Se usa una clave estable, se persiste el resultado junto con el efecto y los reintentos devuelven ese mismo resultado. El costo es estado adicional, política de expiración y coordinación con el almacenamiento.

La comparación debe usar al menos semántica, latencia, disponibilidad, complejidad operativa y recuperación. La alternativa preferible cambia cuando cambia el costo del error en el dominio; no existe ganador universal.

La evidencia apropiada sería curvas de carga, colas, breaker state y degradación del camino crítico. También debe indicarse una condición que refute la elección o haga preferible una solución más simple.
\end{uaanswer}

\begin{uaexercise}
Explique en qué casos retry es razonable y en qué casos es peligroso.
\end{uaexercise}
\begin{uaanswer}
Linearizabilidad exige que cada operación parezca ocurrir en un punto entre invocación y respuesta y que el orden respete tiempo real. Para validar una historia se propone un orden secuencial compatible; si ninguna ubicación de las operaciones satisface ambos requisitos, la historia no es linearizable.

Una solución completa organiza la decisión con P-F-G-S-T: define el problema, enumera fallas, fija la garantía, presenta el protocolo y reconoce el trade-off. Debe incluir estados intermedios y qué ve el usuario si la operación no termina.

La evidencia apropiada sería curvas de carga, colas, breaker state y degradación del camino crítico. También debe indicarse una condición que refute la elección o haga preferible una solución más simple.
\end{uaanswer}

\begin{uaexercise}
Diseñe una política de retry para una operación idempotente.
\end{uaexercise}
\begin{uaanswer}
Una política razonable incluye:
\begin{itemize}
\item máximo de intentos pequeño y explícito;
\item exponential backoff;
\item jitter;
\item timeout acotado por intento;
\item solo aplicar si el error es potencialmente transitorio.
\end{itemize}

La idempotencia es central porque convierte múltiples intentos físicos en un único efecto lógico.
\end{uaanswer}

\begin{uaexercise}
Explique por qué el exponential backoff es mejor que reintentos inmediatos constantes.
\end{uaexercise}
\begin{uaanswer}
La idempotencia separa intentos físicos de una operación lógica. Se usa una clave estable, se persiste el resultado junto con el efecto y los reintentos devuelven ese mismo resultado. El costo es estado adicional, política de expiración y coordinación con el almacenamiento.

El argumento debe identificar la hipótesis, construir la cadena lógica o el contraejemplo y concluir exactamente qué propiedad queda demostrada. Una intuición sin secuencia verificable no es suficiente.

La evidencia apropiada sería curvas de carga, colas, breaker state y degradación del camino crítico. También debe indicarse una condición que refute la elección o haga preferible una solución más simple.
\end{uaanswer}

\begin{uaexercise}
Explique el rol del jitter en un sistema con miles de clientes concurrentes.
\end{uaexercise}
\begin{uaanswer}
La idempotencia separa intentos físicos de una operación lógica. Se usa una clave estable, se persiste el resultado junto con el efecto y los reintentos devuelven ese mismo resultado. El costo es estado adicional, política de expiración y coordinación con el almacenamiento.

La evidencia apropiada sería curvas de carga, colas, breaker state y degradación del camino crítico. También debe indicarse una condición que refute la elección o haga preferible una solución más simple.
\end{uaanswer}

\begin{uaexercise}
Construya un ejemplo donde retries en múltiples capas provoquen una tormenta de requests.
\end{uaexercise}
\begin{uaanswer}
La idempotencia separa intentos físicos de una operación lógica. Se usa una clave estable, se persiste el resultado junto con el efecto y los reintentos devuelven ese mismo resultado. El costo es estado adicional, política de expiración y coordinación con el almacenamiento.

Una solución completa organiza la decisión con P-F-G-S-T: define el problema, enumera fallas, fija la garantía, presenta el protocolo y reconoce el trade-off. Debe incluir estados intermedios y qué ve el usuario si la operación no termina.

La evidencia apropiada sería curvas de carga, colas, breaker state y degradación del camino crítico. También debe indicarse una condición que refute la elección o haga preferible una solución más simple.
\end{uaanswer}

\begin{uaexercise}
Defina circuit breaker y describa sus estados.
\end{uaexercise}
\begin{uaanswer}
Un circuit breaker es un mecanismo que deja de enviar tráfico a una dependencia cuando detecta que hacerlo empeora la situación.

Estados:
\begin{itemize}
\item \textbf{Closed}: tráfico normal;
\item \textbf{Open}: llamadas bloqueadas;
\item \textbf{Half-open}: se dejan pasar algunas llamadas de prueba.
\end{itemize}

La lógica es evitar insistir ciegamente sobre un componente ya degradado.
\end{uaanswer}

\begin{uaexercise}
Explique con un ejemplo qué gana el sistema al abrir un breaker.
\end{uaexercise}
\begin{uaanswer}
La respuesta debe situarse en retry budgets, backoff, jitter, circuit breakers, backpressure, bulkheads y graceful degradation. Se comienza declarando el supuesto decisivo y el comportamiento observable; luego se recorre una ejecución nominal y otra bajo dependencia lenta, retry storm, cola sin límite, pool agotado y cascada. El mecanismo elegido debe vincularse con una garantía concreta, evitando afirmar propiedades fuera de su alcance.

La evidencia apropiada sería curvas de carga, colas, breaker state y degradación del camino crítico. También debe indicarse una condición que refute la elección o haga preferible una solución más simple.
\end{uaanswer}

\begin{uaexercise}
Construya un caso donde un threshold mal calibrado de breaker empeore el comportamiento del sistema.
\end{uaexercise}
\begin{uaanswer}
La idempotencia separa intentos físicos de una operación lógica. Se usa una clave estable, se persiste el resultado junto con el efecto y los reintentos devuelven ese mismo resultado. El costo es estado adicional, política de expiración y coordinación con el almacenamiento.

Una solución completa organiza la decisión con P-F-G-S-T: define el problema, enumera fallas, fija la garantía, presenta el protocolo y reconoce el trade-off. Debe incluir estados intermedios y qué ve el usuario si la operación no termina.

La evidencia apropiada sería curvas de carga, colas, breaker state y degradación del camino crítico. También debe indicarse una condición que refute la elección o haga preferible una solución más simple.
\end{uaanswer}

\begin{uaexercise}
Explique por qué el estado half-open es importante.
\end{uaexercise}
\begin{uaanswer}
Porque evita dos extremos incorrectos:
\begin{itemize}
\item mantener el breaker abierto indefinidamente aunque la dependencia ya se haya recuperado;
\item cerrarlo de golpe y restaurar todo el tráfico sin validación.
\end{itemize}

Half-open permite probar recuperación con tráfico controlado.
\end{uaanswer}

\begin{uaexercise}
Explique qué problema resuelve backpressure.
\end{uaexercise}
\begin{uaanswer}
Backpressure resuelve el problema de desbalance entre ritmo de producción y capacidad de consumo. Si no se regula la entrada, las colas crecen, aumenta la latencia, se consumen buffers y el sistema puede entrar en espiral de degradación.
\end{uaanswer}

\begin{uaexercise}
Diseñe un ejemplo donde la ausencia de backpressure haga crecer una cola hasta volver inútil el sistema.
\end{uaexercise}
\begin{uaanswer}
La idempotencia separa intentos físicos de una operación lógica. Se usa una clave estable, se persiste el resultado junto con el efecto y los reintentos devuelven ese mismo resultado. El costo es estado adicional, política de expiración y coordinación con el almacenamiento.

Una solución completa organiza la decisión con P-F-G-S-T: define el problema, enumera fallas, fija la garantía, presenta el protocolo y reconoce el trade-off. Debe incluir estados intermedios y qué ve el usuario si la operación no termina.

La evidencia apropiada sería curvas de carga, colas, breaker state y degradación del camino crítico. También debe indicarse una condición que refute la elección o haga preferible una solución más simple.
\end{uaanswer}

\begin{uaexercise}
Explique qué es bulkhead y cómo se implementa en términos de recursos.
\end{uaexercise}
\begin{uaanswer}
La idempotencia separa intentos físicos de una operación lógica. Se usa una clave estable, se persiste el resultado junto con el efecto y los reintentos devuelven ese mismo resultado. El costo es estado adicional, política de expiración y coordinación con el almacenamiento.

La evidencia apropiada sería curvas de carga, colas, breaker state y degradación del camino crítico. También debe indicarse una condición que refute la elección o haga preferible una solución más simple.
\end{uaanswer}

\begin{uaexercise}
Construya un caso donde un pool compartido arrastre rutas críticas que no tenían problema.
\end{uaexercise}
\begin{uaanswer}
La idempotencia separa intentos físicos de una operación lógica. Se usa una clave estable, se persiste el resultado junto con el efecto y los reintentos devuelven ese mismo resultado. El costo es estado adicional, política de expiración y coordinación con el almacenamiento.

Una solución completa organiza la decisión con P-F-G-S-T: define el problema, enumera fallas, fija la garantía, presenta el protocolo y reconoce el trade-off. Debe incluir estados intermedios y qué ve el usuario si la operación no termina.

La evidencia apropiada sería curvas de carga, colas, breaker state y degradación del camino crítico. También debe indicarse una condición que refute la elección o haga preferible una solución más simple.
\end{uaanswer}

\begin{uaexercise}
Defina graceful degradation y compárela con caída total.
\end{uaexercise}
\begin{uaanswer}
La idempotencia separa intentos físicos de una operación lógica. Se usa una clave estable, se persiste el resultado junto con el efecto y los reintentos devuelven ese mismo resultado. El costo es estado adicional, política de expiración y coordinación con el almacenamiento.

La comparación debe usar al menos semántica, latencia, disponibilidad, complejidad operativa y recuperación. La alternativa preferible cambia cuando cambia el costo del error en el dominio; no existe ganador universal.

La evidencia apropiada sería curvas de carga, colas, breaker state y degradación del camino crítico. También debe indicarse una condición que refute la elección o haga preferible una solución más simple.
\end{uaanswer}

\begin{uaexercise}
Diseñe tres ejemplos concretos de degradación elegante en una plataforma de e-commerce.
\end{uaexercise}
\begin{uaanswer}
La idempotencia separa intentos físicos de una operación lógica. Se usa una clave estable, se persiste el resultado junto con el efecto y los reintentos devuelven ese mismo resultado. El costo es estado adicional, política de expiración y coordinación con el almacenamiento.

Una solución completa organiza la decisión con P-F-G-S-T: define el problema, enumera fallas, fija la garantía, presenta el protocolo y reconoce el trade-off. Debe incluir estados intermedios y qué ve el usuario si la operación no termina.

La evidencia apropiada sería curvas de carga, colas, breaker state y degradación del camino crítico. También debe indicarse una condición que refute la elección o haga preferible una solución más simple.
\end{uaanswer}

\begin{uaexercise}
Explique qué riesgos aparecen si la degradación no es explícita y semánticamente clara.
\end{uaexercise}
\begin{uaanswer}
La idempotencia separa intentos físicos de una operación lógica. Se usa una clave estable, se persiste el resultado junto con el efecto y los reintentos devuelven ese mismo resultado. El costo es estado adicional, política de expiración y coordinación con el almacenamiento.

La evidencia apropiada sería curvas de carga, colas, breaker state y degradación del camino crítico. También debe indicarse una condición que refute la elección o haga preferible una solución más simple.
\end{uaanswer}

\end{document}
